\documentclass[12pt,letterpaper]{report}
\marginparsep 0pt
\textwidth 6in
\topmargin 0pt
\headsep .5in
\textheight 9.2in
\voffset = 0pt
\hoffset = 0pt
\marginparwidth = 0pt \oddsidemargin = 0pt \sloppy

%Dimensiones de la página
\usepackage[left=2.5cm,top=3cm,right=2.5cm,bottom=2.5cm]{geometry}
%Sangría
\setlength{\parindent}{1cm}

%Numeracion
\pagenumbering{arabic}

\usepackage{template}
\usepackage{amsmath,amsfonts}
\usepackage{comment}
\usepackage{graphicx}
\usepackage{graphics}
\usepackage[dvips]{epsfig}
\usepackage{times}
\usepackage[utf8]{inputenc}
\usepackage[dvips]{graphicx}
\usepackage[usenames]{color}
\usepackage[spanish]{babel}
\newcommand{\ie}{i.e.}
\newcommand {\out}[1]{}
\newtheorem{definicion}{Definicion}
\usepackage[dvips]{epsfig}
\usepackage{rotating}
\usepackage{multirow,multicol}
\usepackage{array}
\usepackage{longtable}
\usepackage[]{fontenc}
\usepackage{hyperref}
\usepackage{float}
\usepackage{ulem}
\usepackage{booktabs}

\usepackage[dvipsnames,svgnames,table]{xcolor}
\usepackage{epstopdf}
\usepackage{color,colortbl,xcolor}
\usepackage{pst-gantt}
\usepackage{ulem}
\usepackage{amssymb}
\setlength{\tabcolsep}{3pt}
\renewcommand{\arraystretch}{1.5}
\usepackage[scaled]{helvet}
\usepackage[shortlabels]{enumitem}
\usepackage{wrapfig}
\usepackage{fancyhdr}
\usepackage{etoolbox}




\usepackage[linesnumbered,ruled,vlined]{algorithm2e}

\usepackage{verbatim}

\hypersetup{
 colorlinks=true,
 linkcolor=black,
 filecolor=black,
 urlcolor=black,
 pdfpagemode=FullScreen,
}

\newcommand{\removed}[1]{{\color{red}\sout{#1}}}
\newcommand{\added}[1]{{\color{blue}#1}}
\newcommand{\modified}[1]{{\color{brown}#1}}
\newcommand{\ra}[1]{\renewcommand{\arraystretch}{#1}}

% Keywords command
\providecommand{\keywords}[1]
{
  \small
  \textbf{\textit{Keywords---}} #1
}

\renewcommand{\shorthandsspanish}{}

\addto\captionsspanish{\def\listtablename{Índice de tablas}
\def\tablename{Tabla}}

\renewcommand{\listfigurename}{Índice de figuras}

\sloppy

\begin{document}
\title{\textbf{Análisis bayesiano con modelos  de XAI sobre métricas de estancamiento en algoritmos bioinspirados para optimización global}\\
%% \title{\textbf{Un enfoque de análisis estadístico de métricas para la mejora estratégica de optimización en algoritmos metaheurísticos}\\
\textit{Bayesian analysis with XAI models on stagnation metrics in bioinspired algorithms for global optimization}}
\author{\textbf{Sebastián M. Guzmán}}
\principaladviser{Rodrigo Olivares O.}
\coprincipaladviser{Nombre Profesor Co-Guía (si aplica)}
\firstreader{Nombre Profesor Informante 1}

\beforepreface


\prefacesection{Resumen}
A continuación aspectos que puede considera para escribir un resumen comprensible y que cumpla su propósito.

\begin{itemize}
    \item Una o dos oraciones que provean una introducción básica al área de trabajo, comprensibles para un interesado de cualquier disciplina.
    \item Dos o tres oraciones de antecedentes más detallados, comprensibles para personas de disciplinas relacionadas.
    \item Una oración que indica claramente el problema general que trata en este trabajo de título.
    \item Una oración que resume el resultado principal.
    \item Dos o tres oraciones que explique lo que el resultado principal aporta al estado del arte.
    \item Una o dos oraciones para poner los resultados en un contexto más general.
    \item (opcional) Dos o tres oraciones para proporcionar una perspectiva más amplia, fácilmente comprensible para una persona de cualquier disciplina.
    \item Máximo 250 palabras.
\end{itemize}

\keywords{uno, dos, tres, cuatro}

\tableofcontents
\newpage

%%=============================================================================
%% List of Figures

\listoffigures
\newpage

%%=============================================================================
%% List of Tables
\renewcommand{\listtablename}{Índice de tablas}
\renewcommand{\tablename}{Tabla}

\listoftables
\newpage

\pagenumbering{arabic}


%==============================
\begin{comment}
 ULTIMA MODIFICACION: 2023
\end{comment}
%==============================

\chapter{Introducción}
   
La investigación en inteligencia artificial (IA) comenzó a ganar popularidad a partir de mediados del siglo XX. El término``inteligencia artificial'' fue propuesto en 1956 durante la conferencia de Dartmouth, que se considera un hito fundacional en el campo de la IA \cite{McCarthy_Minsky_Rochester_Shannon_2006}. Sin embargo, la investigación y el interés en esta área han crecido significativamente en las últimas décadas, especialmente en los años 1980 y 1990 con el avance de las tecnologías de computación y la disponibilidad de mayores capacidades de procesamiento \cite{aiindex2023, LeCun2015}.

La optimización es un elemento crítico en diversas industrias, desde la logística hasta la manufactura, donde se busca maximizar la eficiencia y minimizar costos. Varios métodos tradicionales o heurísticos fueron desarrollados para lograr estos objetivos de optimización \cite{CHEN201771, bard1990branch, vince2002framework}. Más tarde, Fred Glover, acuñó el término ``metaheurística'' en 1986 \cite{glover1986future} para describir métodos heurísticos con características no específicas del problema \cite{abdel2018metaheuristic}. Dentro de la IA, se han destacado subcampos como el aprendizaje automático (Machine Learning), aprendizaje profundo (Deep Learning), y la optimización evolutiva (Evolutionary Algorithms). El aprendizaje automático y profundo abarcan técnicas como las redes neuronales artificiales, máquinas de soporte vectorial, redes convolucionales, entre otras. Estas se utilizan para identificar patrones y realizar predicciones a partir de grandes conjuntos de datos. Recientemente, surgió el término Inteligencia Artificial Explicable (XAI) en respuesta a la creciente complejidad y opacidad de los modelos modernos de IA, especialmente en aplicaciones críticas donde es escencial comprender el proceso de toma de decisiones para la confianza y la responsabilidad. XAI busca hacer que los sistemas de IA sean más transparentes e interpretables.

Por otra parte, la optimización evolutiva utiliza algoritmos basados en principios de evolución natural, como los algoritmos genéticos, para encontrar soluciones óptimas en espacios de búsqueda complejos \cite{izzo2019survey}. En el  contexto de la optimización metaheurística, se han desarrollado diversos algoritmos que proponen técnicas prometedoras para solucionar los problemas de optimización. Algunos de estos algoritmos son bioinspirados como  Ant Colony Optimization (ACO) \cite{484436}, Particle Swarm Optimization (PSO) \cite{488968}, Bat Algorithm (BA) \cite{bat-algorithm}, Dragonfly Algorithm (DA) \cite{Meraihi2020}, Slap Swarm Algorithm (SSA) \cite{MIRJALILI2017163}, entre otros que han surgido en las últimas decadas.

Sin embargo, a pesar de los avances en algoritmos de optimización y sus variantes de mejoras e hibridación, persiste la problemática del estancamiento, donde estos algoritmos se ven atrapados en soluciones subóptimas debido a la falta de exploración en el espacio de búsqueda. Esta problemática afecta a los sectores productivos que utilizan algoritmos de optimización para obtener beneficios diversos. Por ejemplo, las empresas de transporte que deben decidir las rutas más rentables o las organizaciones que enfrentan problemas complejos de optimización. En particular, el Problema del Recorrido Rentable (PTP), que combina beneficios y costos de viaje en una función objetivo, pone de manifiesto la necesidad de estrategias más eficaces para superar el estancamiento y mejorar la toma de decisiones.

Observando en la literatura un fuerte avance en la interpretabilidad de modelos por parte de XAI, la propuesta de este trabajo consiste en un enfoque híbrido que integra modelos bayesianos y técnicas de explicabilidad XAI, como SHAP, para abordar el estancamiento en algoritmos de optimización. Esta metodología no solo busca mejorar la calidad de las soluciones encontradas, sino también proporcionar un entendimiento claro de los factores que contribuyen a la generación de estancamientos.

Las contribuciones de este estudio pueden ser significativas, ya que aportarían un nuevo enfoque para la optimización de algoritmos, abriendo nuevas áreas de investigación en metaheurísticas. Además, la aplicación del método propuesto a un problema real como el PTP ofrecerá no solo soluciones más eficientes, sino que también permitirá a las empresas optimizar sus operaciones y mejorar la satisfacción del cliente. Al combinar la teoría con la práctica, este trabajo busca abordar problemas académicos y ofrecer soluciones aplicables y útiles en el mundo real.


La organización de este trabajo de investigación es la siguiente: el Capítulo 2 presenta el marco conceptual \ref{sc:MC} del trabajo junto al estado del arte \ref{sc:EA}. El Capítulo 3 se expone la descripción del problema \ref{sc:FP}, la propuesta de solución \ref{sc:SP}, preguntas de investigación \ref{sc:PI}, hipótesis \ref{sc:HIPO}, los objetivos generales \ref{sc:OG} y específicos \ref{sc:OE}. En el Capítulo 4 \ref{sc:METOD} se detalla la metodología empleada, incluyendo el diseño de las experimentaciones, el problema de optimización a resolver, el análisis de resultados y cómo serán reportados. El Capítulo 5 describe el plan de trabajo \ref{sc:PT}. Finalmente, en el Capítulo 6 se especifican los recursos humanos \ref{sc:RH} y materiales \ref{sc:RM} requeridos durante la investigación.

\chapter{Marco Teórico y estado del arte}
    
\section{Marco Teórico}
\label{sc:MC}

% ------------------------------------ \\

\subsection{Problemas combinatoriales con restricciones}

\vspace{1mm}
\normalsize

Los problemas de optimización requieren tomar decisiones para seleccionar una solución óptima dentro de un dominio determinado, respetando las restricciones asociadas. El objetivo es cumplir con un criterio específico, que puede ser maximizar o minimizar una determinada medida. Un ejemplo común es maximizar las ganancias de una empresa considerando los costos e ingresos, como en el caso de la venta de camas y veladores. El proceso básico para tomar una decisión en problemas de optimización implica formular el problema, modelarlo, optimizarlo e implementar una solución \cite{10.5555/1718024}.

Para resolver estos problemas mediante algoritmos, se modelan a través de los siguientes componentes:

\begin{itemize}
\item \textbf{Función objetivo f(x)}: es una fórmula matemática que depende de las variables de decisión, coeficientes y una magnitud que representa el objetivo del sistema. El valor resultante de esta función se denota por la letra \textit{z}. La función objetivo busca encontrar el máximo o mínimo de un criterio específico, como maximizar ingresos o minimizar el tiempo de respuesta.

    \item \textbf{Variables de decisión}: son los factores involucrados en la decisión de la función objetivo, representados como $\{x_1, x_2, x_3, ... x_n\}$. Estas variables pueden tener dominios binarios, enteros, reales u otros, dependiendo del problema. Ejemplos incluyen la cantidad de máquinas en una fábrica o la cantidad de computadoras necesarias.

    \item \textbf{Restricciones}: son las limitaciones que deben ser consideradas según el problema. No todos los problemas tienen restricciones, pero cuando existen, pueden incluir recursos, tiempo, etc. Las restricciones se representan mediante ecuaciones o inecuaciones.

\end{itemize}

Las soluciones de estos problemas se representan mediante vectores de tamaño \textit{n}, donde \textit{n} es el número de variables involucradas. Los valores de los vectores solución se sustituyen en la función objetivo para obtener un valor \textit{z}. El conjunto de posibles soluciones que satisfacen las restricciones se denomina \textit{espacio de búsqueda n-dimensional} \cite{10.5555/1718024}. Cada solución es un punto en este espacio, representando una combinación de valores para la función objetivo que se desea maximizar o minimizar \cite{investigacion-operaciones}. La definición formal de un problema de optimización es la siguiente, según \cite{investigacion-operaciones}:

Función objetivo:
\begin{equation}
minimizar/maximizar \quad Z = c_1 x_1 + c_2 x_2 + ... + c_n x_n,
\end{equation}
sujeta a las restricciones:
\begin{align*}
    a_{11} x_1 + a_{12} x_2 + \cdots + a_{1n} x_n \leq b_1 \\
    a_{21} x_1 + a_{22} x_2 + \cdots + a_{2n} x_n \leq b_2 \\
    \cdots \cdots \\
    a_{m1} x_1 + a_{m2} x_2 + \cdots + a_{mn} x_n \leq b_m
\end{align*}
\noindent
y otras restricciones del dominio de $x_i$, donde $i \in\ {1,2, ... ,n}$, $a_{ij}$ es el consumo de recurso por unidad de actividad, $c_{i}$ la contribución a Z por unidad de actividad, y $b_{i}$ la cantidad de recursos disponibles \cite{investigacion-operaciones}.

Las soluciones que cumplen las restricciones se denominan soluciones factibles. Dentro del espacio de búsqueda de soluciones factibles, existen regiones o subconjuntos con un valor \textit{fitness} mejor que otros. Este valor se llama \textit{óptimo local} y puede haber varios en la misma región. La solución con el mejor valor \textit{fitness} para la función objetivo en toda la región factible se llama \textit{óptimo global}.

\vspace{2mm}
%----


\subsection{Heurísticas}

\vspace{1mm}
\normalsize

El concepto de heurística se remonta a la Grecia clásica y, en resumen, significa "encontrar" o "descubrir" algo. En el contexto moderno, la heurística se relaciona con la capacidad de resolver problemas de manera inteligente, aprovechando la información disponible \cite{aprox-heuristica-metaheu}. La historia de la heurística puede caracterizarse como una búsqueda de algoritmos eficientes. Una de las principales motivaciones para utilizar heurísticas es la capacidad de tomar decisiones rápidas y económicas, sin necesariamente comprometer la calidad o la efectividad de la decisión final \cite{heu-tools-for-un-world}.

En el ámbito de las ciencias de la computación, las heurísticas se emplean principalmente para resolver problemas de optimización. Sin embargo, es importante destacar que los métodos y algoritmos de optimización heurística no garantizan encontrar la solución óptima \cite{10.5555/2020934}.

Dentro del campo de la heurística, se han desarrollado varios métodos, la mayoría de los cuales están diseñados para problemas específicos. Algunas de las clasificaciones de heurísticas descritas por \cite{aprox-heuristica-metaheu} incluyen: métodos de descomposición, inductivos, de reducción, constructivos y de búsqueda local. En resumen, todos estos métodos tienen como objetivo generar una solución inicial para un problema y luego, mediante la aplicación de técnicas heurísticas, mejorar gradualmente esta solución hasta alcanzar una aproximación de alta calidad, aunque no necesariamente óptima.

\vspace{2mm}
%----


\subsection{Metaheurísticas (MH)}

\vspace{1mm}
\normalsize

Estas técnicas se consideran una estructura algorítmica estocástica que generalmente se aplica a una variedad de problemas de optimización con solo algunas modificaciones para adaptarse al problema dado. Estas son estrategias para diseñar y/o mejorar los procedimientos heurísticos orientados a obtener un alto rendimiento. Las metaheurísticas proponen un marco general para crear nuevos algoritmos híbridos, combinando varios conceptos, como lo son la genética, biología, inteligencia artificial, matemáticas, física, neurología, entre otras \cite{aprox-heuristica-metaheu}. Estos algoritmos de metaheurística usan alguna compensación de búsqueda local y exploración global aleatoria. Los algoritmos utilizados suelen terminar en un tiempo razonable a como lo haría un algoritmo que pretende buscar dentro de todas las posibles soluciones que se puedan encontrar \cite{metaheu-algo-modeling-optimization}. Por lo que, no hay garantía de encontrar la mejor solución, pero si encontrar una solución de calidad en un tiempo aceptable.

Los principales componentes de los algoritmos de metaheurística son la intensificación y diversificación, también conocidos como explotación y exploración respectivamente. A pesar de que muchos investigadores en sus trabajos han mencionado con frecuencia estos dos conceptos, a menudo, se utilizan definiciones informales y algo vagas de los procesos de estos. Exploración se refiere a generar varias soluciones con el fin de explorar el espacio de búsqueda del problema a escala global, visitando regiones completamente nuevas \cite{metaheu-algo-modeling-optimization}. Mientras que explotación se define como visitar aquellas regiones de un espacio de búsqueda dentro de la vecindad de puntos previamente visitados, centrando la búsqueda dentro del área o región en específica \cite{10.1145/2480741.2480752}. Esto con el fin de que esta última explote la información donde una buena solución fue encontrada en aquella prometedora región.

La diversidad se refiere a las diferencias entre los individuos, es decir, (en términos metaheurísticos) las distancias entre los agentes dentro de la dimensión del espacio de búsqueda en le que se encuentren \cite{10.1145/2480741.2480752}. Según las medidas de diversificación de \cite{10.1145/2480741.2480752}, un aumento en la diversidad indicaba que el algoritmo en cuestión estaba en fase de exploración o ``búsqueda aleatoria'', lo que permite evitar quedar atrapado en óptimos locales de manera temprana, mientras que una disminución en la diversidad indicaba que el procedimiento estaba en fase de explotación o ``búsqueda local''.

En la metaheurística, podemos hablar de que hay dos principales clases. De acuerdo a \cite{10.5555/1718024}, estas son:
\begin{itemize}
    \item \textbf{Single-solution-based} o \textbf{Trajectory-based}: el concepto principal detrás de este, es aplicar iterativamente los procedimientos de generación y reemplazo de la solución única actual. El proceso se comienza, en la fase de generación. Donde un conjunto de soluciones candidatas se generan a partir de la solución actual $s$. Este conjunto $C(s)$ se obtiene generalmente por transformaciones locales de la solución. En la fase de reemplazo, se realiza una selección del conjunto de soluciones candidatas $C(s)$ para reemplazar la solución actual; es decir, se selecciona una solución $s' \in C(s)$ para que sea la nueva solución. Este proceso itera hasta un criterio de parada entregado previamente \cite{10.5555/1718024}.

    \item \textbf{Population-based}: Las metaheurísticas basadas en esta categoría parten de una población inicial de soluciones. Luego, iterativamente aplican la generación de una nueva población y el reemplazo de la población actual. En la fase de generación, se crea una nueva población de soluciones. En la fase de reposición se realiza una selección entre las poblaciones actuales y las nuevas. Este proceso itera hasta un criterio de parada dado \cite{10.5555/1718024}.

\end{itemize}

Dentro de la clase de metaheurísticas population-based, existen algoritmos bio-inspirados (de la biología), uno de ellos es el interesante paradigma llamado \textit{Swarm Intelligence (SI)}. Los algoritmos basados en el comportamiento colectivo de especies como lobos, hormigas, abejas, avispas, termitas, peces, pájaros, etcetera, se denominan algoritmos SI \cite{10.5555/1718024}. Este tipo de algoritmos se originaron a partir del comportamiento social de aquellas especies que compiten por los alimentos. Las características principales de los algoritmos SI son que las partículas son agentes simples que cooperan mediante una comunicación indirecta, y que realizan movimientos en el espacio de decisión \cite{10.5555/1718024}.

\vspace{2mm}
%-----


\subsection{Óptimo global y óptimo local}

\vspace{1mm}
\normalsize
En el contexto de las metaheurísticas, los conceptos de óptimo global y óptimo local son fundamentales para entender su funcionamiento y desafíos \cite{10.5555/1718024}.
\begin{itemize}
    \item \textbf{Óptimo global}: es la mejor solución posible para un problema de optimización, la que tiene el valor más favorable de la función objetivo en todo el espacio de búsqueda.
    \item \textbf{Óptimo local}: es una solución que es mejor que todas las soluciones en su vecindad inmediata, pero no necesariamente la mejor en todo el espacio de búsqueda .
\end{itemize}



\begin{figure}[ht!]
    \centering
    \includegraphics[width=0.6\textwidth]{imagenes/localAndGlobalOptima.png}
    \caption{Imagen de ejemplo de óptimo local y global con función de minimización.}
    \label{fig:ex_opt_local_global}
\end{figure}

En la Figura \ref{fig:ex_opt_local_global} se puede ver una función objetivo que toma una variable unidimensional, la cual es evaluada y toma un valor. La función es de minimización, por lo que un algoritmo podría encontrar el mínimo valor de una curva en zonas convexas en esta función varias veces, que podría no ser el óptimo global (el punto rojo), estos son los puntos azules, los óptimos locales.







\vspace{2mm}
%-----



\subsection{Análisis Estadístico de Variables y Modelos de Clasificación}

\vspace{1mm}
\normalsize

El análisis estadístico de variables es una parte fundamental de la ciencia de datos que implica el estudio y la comprensión de las relaciones entre variables, así como la construcción de modelos predictivos para clasificar datos en diferentes categorías. En este contexto existen dos ramas de la estadística, la descriptiva y la bayesiana. La estadística descriptiva se enfoca en describir y resumir características clave de un conjunto de datos, como la media, la mediana y la desviación estándar, proporcionando una comprensión clara de la distribución de los datos. Por otro lado, la estadística bayesiana utiliza el teorema de Bayes para actualizar las creencias sobre los parámetros del modelo a medida que se observan datos adicionales, permitiendo una estimación más completa de la incertidumbre en las inferencias y una toma de decisiones más informada en base a la evidencia observada y tomar decisiones basadas en esta información \cite{McClave2016_lz}.


\subsubsection{Estadística Bayesiana}

Como se dijo, la estadística bayesiana es un enfoque estadístico que se basa en el teorema de Bayes para actualizar las creencias sobre los parámetros de un modelo a medida que se observan datos adicionales. En lugar de tratar los parámetros como valores fijos pero desconocidos, la estadística bayesiana los modela como variables aleatorias con distribuciones de probabilidad. Al combinar la información previa (a priori) con la evidencia observada (verosimilitud), se obtiene una distribución posterior que representa la incertidumbre sobre los parámetros después de observar los datos. Esto permite una estimación más completa de la incertidumbre en las inferencias y una toma de decisiones más informada \cite{Gelman2013_ay}.

El teorema de Bayes se expresa como:

\begin{equation}
P(\theta|D) = \frac{P(D|\theta) \cdot P(\theta)}{P(D)}
\end{equation}

Donde:
\begin{itemize}
    \item $P(\theta|D)$ es la distribución posterior de los parámetros $\theta$ dado el conjunto de datos $D$.
    \item $P(D|\theta)$ es la verosimilitud del conjunto de datos $D$ dado el parámetro $\theta$.
    \item $P(\theta)$ es la distribución a priori de los parámetros $\theta$.
    \item $P(D)$ es la verosimilitud marginal o evidencia.
\end{itemize}


\subsubsection{Modelos de Clasificación}

Los modelos de clasificación son herramientas estadísticas que se utilizan para asignar una categoría o etiqueta a una observación en función de sus características. En el contexto de la estadística bayesiana, los modelos de clasificación bayesianos utilizan la probabilidad condicional para calcular la probabilidad de pertenencia a cada clase dado un conjunto de características observadas. Estos modelos pueden ser simples, como el clasificador bayesiano ingenuo (Naive Bayes), que asume independencia condicional entre las características, o más complejos, como los modelos de regresión logística bayesianos o las máquinas de vectores de soporte bayesianas \cite{Murphy2012_nx}.

Estos modelos de clasificación bayesianos no solo proporcionan predicciones de clase, sino también estimaciones de incertidumbre que pueden ser útiles en la toma de decisiones. Además, la estadística bayesiana permite la inclusión de información previa sobre los parámetros del modelo, lo que puede ser especialmente útil cuando se tienen pocos datos o se requiere una actualización continua del modelo a medida que se obtienen más observaciones.

\vspace{2mm}
%-----




\subsection{Surrogate Fitness Models}

\vspace{1mm}
\normalsize

Los Surrogate Fitness Models, también conocidos como modelos de fitness sustitutos, son herramientas útiles en la optimización con metaheurísticas, ya que permiten predecir los valores de fitness a partir de la función objetivo sin necesidad de evaluarla directamente en cada iteración. Estos modelos se construyen utilizando técnicas de aprendizaje supervisado, como la regresión lineal, polinómica, las máquinas de soporte vectorial (SVM) y las redes neuronales. A partir de un conjunto inicial de datos de entrenamiento generado en las primeras etapas del proceso de optimización, se ajusta un modelo que puede estimar con precisión los valores de fitness de nuevas soluciones. Esto resulta en una evaluación más rápida y eficiente, crucial en problemas de optimización donde las evaluaciones de la función objetivo son costosas en términos de tiempo y recursos computacionales \cite{10138291, Jin2005_te}.


\vspace{2mm}
%-----





\subsection{eXplainable Artificial Intelligence (XAI)}

\vspace{1mm}
\normalsize

La explicabilidad en el aprendizaje automático o XAI es una disciplina que busca hacer comprensibles los modelos de inteligencia artificial para los humanos. Este concepto es crucial para fomentar la confianza, facilitar la toma de decisiones y garantizar la transparencia. Existen varios enfoques para lograr la explicabilidad según \cite{das2020opportunitieschallengesexplainableartificial}, los cuales se pueden categorizar en base a diferentes criterios:

\subsubsection{Enfoques de XAI}

\begin{enumerate}
    \item{Explicabilidad intrínseca vs. explicabilidad post-hoc}
    \begin{itemize}
        \item \textbf{Intrínseca}: Modelos que son interpretables por naturaleza debido a su simplicidad y estructura, como los árboles de decisión y las regresiones lineales.
        \item \textbf{Post-hoc}: Métodos aplicados a modelos complejos (como redes neuronales o ensamblajes de modelos) para explicar sus decisiones después de que han sido entrenados.
    \end{itemize}
\end{enumerate}


\begin{enumerate}
    \item{Explicabilidad global vs. Local}
    \begin{itemize}
        \item \textbf{Global}: Ofrece una comprensión de cómo funciona el modelo en su totalidad.
        \item \textbf{Local}: Se centra en explicar una única predicción o un pequeño conjunto de predicciones.
    \end{itemize}
\end{enumerate}


\begin{enumerate}
    \item{Métodos agnósticos vs. Específicos de modelo}
    \begin{itemize}
        \item \textbf{Agnósticos}: Pueden ser aplicados a cualquier tipo de modelo, como LIME (Local Interpretable Model-agnostic Explanations) y SHAP (SHapley Additive exPlanations) \cite{ribeiro2016why}.
        \item \textbf{Específicos de Modelo}: Diseñados para explicar modelos específicos, como los gráficos de importancia de características en árboles de decisión.
    \end{itemize}
\end{enumerate}


\begin{enumerate}
    \item{Técnicas de explicabilidad comunes}
    \begin{itemize}
        \item \textbf{Importancia de Características}: Mide la contribución de cada característica a las predicciones del modelo.
        \item \textbf{Visualización de Gradientes}: Usada principalmente en redes neuronales para visualizar cómo los cambios en las entradas afectan las salidas.
        \item \textbf{Métodos de Descomposición}: Descomponen la predicción en componentes aditivos que representan la contribución de cada característica, como SHAP.
    \end{itemize}
\end{enumerate}

\vspace{2mm}
%-----



\subsection{SHapley Additive exPlanations (SHAP)}

\vspace{1mm}
\normalsize

SHAP (SHapley Additive exPlanations) es un método post-hoc y agnóstico que proporciona explicaciones locales y globales de las predicciones de los modelos de aprendizaje automático. Basado en la teoría del valor de Shapley de la teoría de juegos, SHAP asigna un valor a cada característica que refleja su contribución a la predicción del modelo. Los valores de Shapley consideran todas las posibles combinaciones de características, proporcionando una medida justa y equitativa de la importancia de cada característica \cite{lundberg2017unifiedapproachinterpretingmodel}.

El enfoque SHAP se destaca por varias razones:
\begin{itemize}
    \item \textbf{Equidad y Consistencia}: Utiliza la teoría de juegos para garantizar que las contribuciones se calculen de manera justa y consistente.
    \item \textbf{Interpretabilidad Global y Local}: Puede descomponer la predicción de un modelo en contribuciones individuales de características para una sola predicción (explicabilidad local) y sumar estas contribuciones para obtener una comprensión global.
    \item \textbf{Compatibilidad Agnóstica}: Se puede aplicar a cualquier tipo de modelo, lo que lo hace extremadamente versátil.
\end{itemize}

El proceso de SHAP implica calcular el valor de cada característica como la media ponderada de las diferencias en las predicciones que resultan de añadir esa característica a todas las combinaciones posibles de otras características. Esta metodología no solo permite comprender cómo cada característica individual influye en la predicción, sino que también facilita la detección de interacciones entre características.

\vspace{2mm}
%-----








\section{Estado del Arte}
\label{sc:EA}

Esta sección recopila trabajos relacionados que han abordado el problema de estancamiento de soluciones en algoritmos de optimización, donde varios enfoques han sido propuestos en distintos algoritmos para resolver el problema del estancamiento de soluciones, que se encuentra estrechamente relacionado al balance de exploración y explotación que realiza el algoritmo durante el proceso de optimización.

\textbf{Enfoques Basados en Inteligencia Artificial Mediante Técnicas Bayesianas:}

En \cite{ZHANG2015138}, se presenta un enfoque avanzado para el algoritmo particle swarm optimizer (PSO) mediante la incorporación de técnicas bayesianas para el ajuste adaptativo del parámetro peso de inercia (que es propio de este algoritmo) que mejora significativamente el rendimiento del algoritmo a través del equilibrio entre exploración y explotación. El algoritmo en \cite{10.1371/journal.pone.0048710}, presenta una interpretación bayesiana del PSO, proporcionando un marco formal para incorporar conocimiento previo sobre el problema a resolver. Lo que hace, es extender el método de optimización mediante funciones kernel, que transforman los datos a un espacio diferente donde el problema puede ser más fácil de resolver. Esta transformación se considera una forma de conocimiento previo sobre la naturaleza del problema de optimización. Este enfoque se aplica para resolver problemas de optimización donde la evaluación de la función de desempeño es costosa en términos de tiempo. En \cite{logistics8010014}, se formula una técnica novedosa que combina la optimización bayesiana con los algoritmos genéticos de agrupamiento. Esta técnica utiliza un enfoque basado en la optimización bayesiana para ajustar los parámetros del Grouping Genetic Algorithms (GGA) de manera más eficiente y eficaz, permitiendo la capacidad para encontrar buenos valores de parámetros con menos evaluaciones de función en comparación con métodos tradicionales. La diferencia del enfoque bayesiano adoptado en este trabajo es que los parámetros que ajustan son tasa de cruce, tasa de mutación, tamaño relativo del grupo de apareamiento, proporción de la élite dentro de la población, probabilidades de aplicar variantes de cruce y otros vistos en la investigación.

\textbf{Enfoques Basados en Combinación Híbrida de Operadores Heurísticos:}

En \cite{wang2021hybrid} se aborda la ineficacia del Butterfly Optimization Algorithm (BOA) en mantener la diversidad y evitar la convergencia prematura en problemas de optimización global que es propenso a caer en óptimos locales y puede perder su capacidad exploratoria en las etapas iniciales de búsqueda. Este trabajo combina BOA y el Flower Pollination Algorithm (FPA), basado en un mecanismo de mutualismo que fundamenta en la capacidad exploratoria del FPA y la capacidad explotadora del BOA, donde las mariposas y las flores se benefician mutuamente: las mariposas obtienen néctar y hábitat, mientras que las flores aseguran la polinización. Otro trabajo \cite{8962051}, trabaja con el algoritmo Ant Colony Optimizer (ACO) que tradicionalmente ha enfrentado problemas al resolver el Travel Salesman Problem (TSP), tales como caer en óptimos locales y tener una convergencia lenta. En por esto que han propuesto un nuevo algoritmo de colonia de hormigas multi-colonia que utiliza Redes Generativas Adversarias (GAN) y una estrategia adaptativa de evitación de estancamiento (GAACO). Enfoque que hizo mejorar la velocidad de convergencia y evitar la convergencia prematura. En \cite{a_hybrid_WOA_seagullAlgorithm}, se propone un nuevo algoritmo híbrido llamado Whale Optimization with Seagull Algorithm (WSOA). Este algoritmo combina el mecanismo de contracción circundante del Whale Optimization Algorithm (WOA) con el comportamiento de ataque en espiral del Seagull Optimization Algorithm (SOA). Además, se introduce la estrategia de vuelo Lévy en la fórmula de búsqueda del SOA para evitar la convergencia prematura y equilibrar de manera más efectiva la exploración y explotación del algoritmo. En \cite{ccelik2021advancement}, trabaja con el algoritmo Slap Swarm Algorithm (SSA), el cual es efectivo en varios casos, pero presenta problemas como la incapacidad de intercambiar información entre los líderes del enjambre y una convergencia prematura, lo que reduce la precisión de la solución y la capacidad de exploración del algoritmo. Este estudio modifica SSA con un mapa sinusoidal para modificar caóticamente el parámetro que equilibra la exploración y explotación, se añade una relación mutualista entre dos líderes del enjambre para mejorar el intercambio de información, y se aplica una técnica aleatoria a los seguidores del enjambre para introducir diversidad y mejorar la exploración de áreas no visitadas del espacio de búsqueda.

\textbf{Enfoques Basados en Optimización Bioinspirada y Algoritmos Evolutivos:}

En \cite{ABDELMAGEED2022107904}, se propone un algoritmo denominado Improved Binary Adaptive Wind Driven Optimization (iBAWDO). Este algoritmo combina técnicas de optimización inspiradas en la naturaleza con operadores evolutivos (como el cruce, la mutación y la hibridación con el algoritmo Simulated Annealing, SA) para mejorar la capacidad de exploración y explotación del espacio de soluciones. Principalmente, se aborda el problema de la selección de características en la clasificación supervisada, donde la alta dimensionalidad de los datos puede afectar negativamente el rendimiento de los clasificadores. Sin embargo, este algoritmo mejorado resulta ser efectivo durante el proceso de optimización. En \cite{garcia2020enhancing}, se propone mejorar un marco de binarización existente que utiliza la técnica K-means incorporando un operador de perturbación basado en la técnica K-nearest neighbor (KNN). Este operador de perturbación tiene como objetivo aumentar la diversificación y la intensificación de los algoritmos metaheurísticos en su versión binaria. En el estudio se utilizaron técnicas como Particle Swarm Optimization (PSO) y Cuckoo Search (CS) para demostrar la eficacia del marco mejorado. En \cite{seyyedabbasi2021gwo}, también trata de resolver el problema de la convergencia temprana en óptimos locales, pero esta vez en el algoritmo Grey Wolf Optimizer (GWO). En este trabajo se introducen dos algoritmos mejorados basados en el anterior, Expanded GWO e Incremental GWO. En el primero se adopta una estrategia que permite que los lobos actualicen sus posiciones en función de los tres mejores lobos (alfa, beta y delta) así como de sus propias posiciones previas y en el segundo enfoque cada lobo actualiza su posición basándose en todos los lobos seleccionados antes que él, permitiendo una mayor diversidad en la búsqueda de soluciones. Esto con el objetivo de mejorar la capacidad de exploración y explotación para encontrar soluciones más eficaces en problemas con múltiples óptimos locales. En \cite{rauf2020particle}, aborda las limitaciones en la inicialización de la población en los algoritmos de optimización por enjambre de partículas (PSO). Para resolver esta problemática el estudio propone una nueva técnica de inicialización de la población utilizando secuencias de probabilidad, específicamente una distribución Weibull. Esta técnica, denominada WI-PSO (Weibull Initialized PSO), aplica la distribución de probabilidad para generar números en ubicaciones aleatorias para la inicialización del enjambre y así  mejorar la diversidad del enjambre y la velocidad de convergencia en la optimización global. La propuesta es comparada con el mismo algoritmo PSO, pero que ha sido utilizado con otras probabilidades para la inicialización de la población. En \cite{zhang2020modified}, se propusieron seis estrategias diferentes para actualizar la energía de escape (E) en el Harris Hawks Optimizer (HHO). Estas estrategias fueron diseñadas para modelar mejor las situaciones reales de caza de los Harris Hawks en la naturaleza. Mediante un estudio comparativo utilizando veinte funciones de referencia representativas, se evaluaron estas estrategias para determinar la más efectiva. La primera estrategia trata de la aplicación de decreciente lineal a \textit{E}, la segunda y tercera de decrecientes no lineales basadas en poder, la cuarta una estrategia convexa-cóncava de una función seno de traslación, la quinta seno cóncavoconvexa (variante de seno), y sexta, con función exponencial. En donde se elige mediante experimentos la mejor de estas, y luego la mejor se compara con otros algoritmos de optimización conocidos.


\textbf{Enfoques Basados en Metaheurísticas Explicables (XAI):}

En \cite{wallace2021towards}, aborda el problema de la falta de justificación en las soluciones generadas por metaheurísticas, las cuales son algoritmos de búsqueda aleatorizados utilizados para encontrar soluciones "suficientemente buenas" a problemas de optimización. El trabajo se enfoca en identificar las combinaciones de variables que influyen fuertemente en la calidad de la solución y la naturaleza de esta relación, lo que representa un paso hacia la explicación de las decisiones tomadas por una metaheurística. La propuesta del artículo consiste en utilizar funciones de fitness sustitutas (surrogate fitness functions) dentro de una metaheurística para identificar las variables importantes en la calidad de la solución generada. Estas funciones sustitutas son modelos matemáticos que se ajustan a las evaluaciones reales de la función objetivo, permitiendo la evaluación rápida de la calidad de la solución y la interpretación del impacto de cada variable en la solución final. En \cite{towardsXAImh_featureExtraction}, se propone un enfoque para la extracción de características explicativas de las metaheurísticas mediante la minería de trayectorias (trajectory mining). La minería de trayectorias se refiere al proceso de análisis y extracción de patrones a partir de las trayectorias seguidas por las soluciones durante el proceso de optimización. En el caso de las metaheurísticas, las trayectorias representan las secuencias de movimientos y cambios en las soluciones generadas a lo largo de las iteraciones del algoritmo. Este enfoque incluye el uso de análisis de componentes principales (Principal Component Analysis, PCA) para reducir la dimensionalidad de los datos de las trayectorias y facilitar la identificación de patrones significativos. Posteriormente, se aplica una técnica de clustering para agrupar las trayectorias en categorías distintas, lo que permite la interpretación de las características de cada grupo y su influencia en el rendimiento de la metaheurística. Este enfoque permite a los investigadores y usuarios de metaheurísticas entender mejor cómo los cambios en las soluciones afectan el desempeño del algoritmo y proporciona una base para mejorar la explicación y la interpretabilidad de los resultados obtenidos. En \cite{singh2022towards}, se utiliza una técnica de optimización basada en modelos de ajuste sustitutos (surrogate modeling) que sirve como proxy para la metaheurística original. En el contexto de las metaheurísticas, un modelo de ajuste sustituto es una aproximación matemática que emula el comportamiento de la metaheurística en términos de sus decisiones y resultados generados. Este enfoque permite descomponer el proceso de búsqueda en una serie de pasos explicables, identificando las interacciones y combinaciones de variables que tienen un impacto significativo en la calidad de la solución. Los modelos de ajuste sustitutos pueden ser de varios tipos, como regresión lineal, árboles de decisión, redes neuronales, entre otros. La técnica propuesta consiste en entrenar un modelo sustituto utilizando datos generados por la metaheurística original durante su proceso de búsqueda. Luego, este modelo se utiliza para analizar y explicar las decisiones tomadas por la metaheurística en términos de las variables de entrada y su influencia en la calidad de la solución. En el contexto de los Algoritmos Genéticos (GA), la técnica se aplica para descomponer y entender los factores que contribuyen a la calidad de las soluciones generadas por el algoritmo. Esta técnica ofrece una manera de interpretar y explicar las decisiones tomadas por una metaheurística, lo que puede ser útil tanto para mejorar el diseño del algoritmo como para aumentar la confianza de los usuarios en los resultados obtenidos.

% tengo justo 15 trabajos relacionados

\chapter{Definición del problema}
   

\section{Formulación del Problema}
\label{sc:FP}

Habiendo revisado varios trabajos en la literatura, se observa que continúa la investigación de diversos métodos para mejorar los algoritmos de optimización con diferentes enfoques. A pesar de los avances, siempre se enfrenta el problema del estancamiento en las metaheurísticas.


\subsection{El Problema del Estancamiento en Metaheurísticas}

El estancamiento en las metaheurísticas es un fenómeno que ocurre cuando un algoritmo de optimización, que busca encontrar la solución óptima o una solución cercana al óptimo en un espacio de búsqueda, se queda atrapado en una región del espacio de soluciones y es incapaz de encontrar mejores soluciones a pesar de seguir iterando, quedando estancado en lo que se llama un óptimo local. Este problema se presenta comúnmente en los algoritmos bio-inspirados que resuelven problemas de optimización combinatorial, donde el espacio de búsqueda es vasto y complejo \cite{10.5555/1718024, Michalewicz2004}.

\subsubsection{Características del Estancamiento}

El estancamiento se caracteriza por:
\begin{itemize}
    \item \textbf{Convergencia Prematura:} El algoritmo converge rápidamente a una solución subóptima y no explora adecuadamente otras regiones del espacio de búsqueda.
    \item \textbf{Falta de Diversidad:} La población de soluciones en algoritmos basados en población (como los algoritmos genéticos) se vuelve homogénea, reduciendo la capacidad de encontrar nuevas soluciones.
    \item \textbf{Exploración Insuficiente:} El algoritmo explora insuficientemente el espacio de búsqueda, centrando sus esfuerzos en áreas ya conocidas y dejando de lado regiones potencialmente prometedoras.
\end{itemize}

\subsubsection{Causas del Estancamiento}

Según \cite{Dokeroglu2019_jx} y de acuerdo a lo revisado en la literatura, las principales causas del estancamiento incluyen:
\begin{itemize}
    \item \textbf{Función de Evaluación Engañosa:} La función objetivo puede tener múltiples óptimos locales que engañan al algoritmo para que se quede atrapado en una de estas soluciones.
    \item \textbf{Parámetros del Algoritmo:} Parámetros mal ajustados, como tasas de mutación y cruce en algoritmos genéticos u otros, pueden llevar a una exploración insuficiente del espacio de búsqueda.
    \item \textbf{Estrategia de Búsqueda:} Estrategias de búsqueda que favorecen la explotación sobre la exploración pueden aumentar la probabilidad de estancamiento.
\end{itemize}

\subsubsection{Consecuencias del Estancamiento}

El estancamiento puede llevar a:
\begin{itemize}
    \item \textbf{Soluciones Subóptimas:} El algoritmo puede terminar con soluciones que están lejos del óptimo global.
    \item \textbf{Tiempo de Cómputo Elevado:} Más iteraciones sin mejora significativa resultan en un uso ineficiente de recursos computacionales.
\end{itemize}

\subsubsection{Métodos para Superar el Estancamiento}

Para superar el estancamiento, en el estado del arte se pudo encontrar varias maneras, y se pueden aplicar diversas estrategias como:
\begin{itemize}
    \item \textbf{Métodos Híbridos:} Combinar diferentes algoritmos de optimización para aprovechar sus fortalezas.
    \item \textbf{Técnicas de Diversificación:} Introducir mecanismos que aumenten la diversidad de la población de soluciones.
    \item \textbf{Reinicialización:} Reiniciar el algoritmo desde una nueva posición en el espacio de búsqueda.
    \item \textbf{Ajuste Dinámico de Parámetros:} Modificar los parámetros del algoritmo durante la ejecución para mejorar la exploración y explotación.
\end{itemize}

Varios de estos métodos utilizan distintos enfoques que manipulan los parámetros de las metaheurísticas. En este trabajo se propone otra manera de resolver el problema del estancamiento utilizando un modelo estadístico bayesiano que pueda ser explicado por SHAP para salir del estancamiento ingeniosamente según las métricas que nos entregue el algoritmo en cuestión.










\section{Solución Propuesta}
\label{sc:SP}

Considerando la falta de trabajos realizados en el contexto de metaheurísticas que sean explicables, en este estudio se propone un método híbrido hacia algoritmos de optimización que les permita utilizar las ventajas de entrenar a un modelo bayesiano durante el proceso de optimización (online), como un modelo que predice el fitness, en conjunto a un método de detección del estancamiento de las soluciones, y que mediante herramientas de explicabilidad como SHAP, se pueda determinar cuál esta siendo el factor causante del estancamiento a través del modelo bayesiano, lo que permitirá actualizar las soluciones de una manera ingeniosa y enfocada al problema.


\begin{figure}[ht!]
    \centering
    \includegraphics[width=0.6\textwidth]{imagenes/diagramaFlujoPropusta.png}
    \caption{Diagrama de flujo de la propuesta de solución.}
    \label{fig:dflujo_propuesta}
\end{figure}

La propuesta se ve en la Figura \ref{fig:dflujo_propuesta}. Se presenta la estructura general de una metaheurística, la que abarca varias implementaciones de algoritmos bio-inspirados. Se comienza con la inicialización de las variables, la población inicial de soluciones y el fitness de estas.

Para esta propuesta, cada uno de los individuos de la población tendrá un Modelo Bayesiano Regresivo (MBR), al cual se le entregará el vector solución y el fitness calculado por cada iteración, a modo de que se entrene de manera "online", a medida que datos nuevos van surgiendo. Esto es lo que representa el color verde del diagrama.

Por otra parte, el algoritmo poseerá un método de detección de estancamiento, plasmada en color naranjo dentro de la figura. En especial, para este estudio se utilizará el ``Monitoreo del cálculo de la Mejora Relativa'' (MMR). Este funcionará de manera individual, es decir, irá tomando el historial de los mejores fitness que se han actualizado de este agente para determinar después de la 1/4 parte del criterio de parada (ventana de detección) si el individuo se encuentra estancado o no. Para calcular la MMR se utiliza la siguiente fórmula:

La mejora relativa se define como:

\[
\text{Mejora Relativa} = \frac{f_1 - f_n}{f_1}
\label{ec:mejora_relativa_1}
\]

La condición para determinar si ha habido una mejora significativa es:

\[
\frac{f_1 - f_n}{f_1} > \epsilon
\label{ec:mejora_relativa_2}
\]

donde \( f_1 \) es el valor de fitness al inicio de la ventana de detección, \( f_n \) es el valor de fitness en la iteración más reciente dentro de esa ventana, y \(\epsilon\) es el umbral de mejora, que preliminarmente se le asignará un valor de $0.01$.

Por último, los procesos de color azul son la parte de la sección de explicabilidad (XAI) del comportamiento de las variables de decisión con respecto al fitness. Este proceso será utilizado solo en caso de que el individuo haya sido detectado en estancamiento. Para explicar las variables que están influyendo se utiliza el MBR propio del agente, que posee el historial de datos que ha utilizado para ir actualizando el fitness. La técnica empleada en esta sección es SHAP, que permitirá ver cuales son las variables que menos están aportando, con el fin de actualizar al individuo con las variables que realmente están aportando más cambio al fitness (maximización o minimización).

El método para salir del estancamiento, se realiza con la selección del número de variables que serán tomadas en cuenta de SHAP, que dependerán de la probabilidad de crossover o mutación que exista, y principalmente del porcentaje de importancia que tenga unas variables respecto a otras.

Una vez, se ejecuta el método para salir del estancamiento, se reentrena el MBR, se comparan las soluciones, se actualiza la mejor solución global de la población, y se verifica el criterio de parada para continuar con el ciclo o finalizarlo. El pseudocódigo de este diagrama puede ser visto en detalle \ref{alg:metaheuristics_bayesianXAI}.



\begin{algorithm}[!ht]
    \scriptsize
    \caption{Estructura común de las metaheurísticas con la propuesta en peudocódigo}\label{alg:metaheuristics_bayesianXAI}
    \KwIn{datos de entrada del problema, $popSize$, $maxIter$, $detectionWindow$, $improvement_threshold$}
    \KwOut{una solución óptima obtenida durante el proceso de optimización.}

    cargarParámetrosDelAlgoritmo()\;
    cargarDatosDelProblema()\;
    \tcp{el valor $n$ define el número de dimensiones}
    Define la función objetivo $f(x)$, $x=\langle x^{1},\ldots,x^{n} \rangle$\;
    \tcp{produce la primera generación de $popSize$ agentes}
    \ForEach{agente $a_{i}$ \textbf{desde} 1 \textbf{hasta} $popSize$}{
        \ForEach{dimensión $j$ \textbf{desde} 1 \textbf{hasta} $n$}{
            \While{posición $x_{i}^{j}$  no sea factible}{
                posición $x_{i}^{j} \leftarrow$ ValorAleatorioEntre\{0, 1\}\;
                Inicialización de variables adicionales de la metaheurística\;
            }
        }
        Maximizar o minimizar la función objetivo y comparar para obtener la mejor solución inicial\;
        $MBR\_x_{i}$ $\leftarrow$ Entrenar Modelo Bayesiano Regresivo del agente $i$\;
    }
    \tcp{producir $maxIter$-generaciones de $popSize$ agentes}
    \While{$t < maxIter$}{
        \tcp{fase del proceso es invocado}
        Asigna valores de diversificación o intensificación (dependiendo de $t$)\;
        \ForEach{agente $a_{i}$ \textbf{desde} 1 \textbf{hasta} $popSize$}{
            \ForEach{dimensión $j$ \textbf{desde} 1 \textbf{hasta} $n$}{
                \tcp{genera nuevas soluciones}
                \While{posición $x_{i}^{j}$ no es factible}{
                    posición $x_{i}^{j} \leftarrow$ es actualizada por las ecuaciones de la metaheuristica\;
                }
            }
            Maximizar o minimizar la función objetivo (calculo del fitness)\;

            actualizar $MBR\_x_{i}$ entrenándolo en línea con nueva data del agente $i$\;

            \If{($t > detectionWindow$)}{
                MMR $\leftarrow$ calcular método de monitoreo de la mejora relativa de la Ec. \ref{ec:mejora_relativa_1} y \ref{ec:mejora_relativa_2}\;
                \If{$MMR$ < $improvement_threshold$}{
                    ejecutar método XAI a través de SHAP para explicar $MBR$ de la solución $i$\;
                    $arrayResponsable$ $\leftarrow$ seleccionar alguna(s) variable(s) responsable(s) estancamiento\;
                    \If{randomValorEntre(0,1) > $probMutacion$}{
                        Ejecutar mutación en $arrayResponsable$\;
                    }
                    \Else{
                        \If{randomValorEntre(0,1) > $probCrossover$}{
                            Ejecutar cruce (crossover) en $arrayResponsable$\;
                        }
                    }
                    Maximizar o minimizar la función objetivo (calculo del fitness)\;

                    actualizar $MBR\_x_{i}$ entrenándolo en línea con nueva data del agente $i$\;
                }
            }
            actualizar la mejor solución \;
        }
    }
    Retornar resultados y visualización post-proceso\;
\end{algorithm}

La propuesta de mejora a las metaheurísticas en el ámbito académico, abre nuevas áreas de investigación y aplicación, fomentando la innovación y el desarrollo de nuevas técnicas y algoritmos. Esto puede llevar a la formación de nuevas teorías y modelos, enriqueciendo el conocimiento en el campo de la optimización. Además, puede facilitar la resolución de problemas complejos de la vida real de manera ingeniosa en diferentes disciplinas.

Además puede tener un impacto significativo en varios ámbitos si es que se aplica en la problemas reales. En el ámbito comercial, empresas que dependen de la optimización de recursos y procesos, como las de logística y manufactura, pueden beneficiarse enormemente de soluciones más eficientes. Esto podría resultar en una reducción de costos operativos, mejor utilización de recursos y una mayor competitividad en el mercado. En el ámbito social, mejoras en la eficiencia de procesos industriales y de transporte pueden contribuir a una reducción de la huella de carbono, promoviendo prácticas más sostenibles y amigables con el medio ambiente.







\section{Pregunta de Investigación}
\label{sc:PI}

Basándonos en la revisión de la literatura, se plantean las siguientes preguntas de investigación:

%Desarrolla al menos una pregunta de investigación, claramente relevante para la validación del tratamiento. (estas podrán ser respondidas luego de haber ejecutado los casos experimentales definidos)%

\begin{itemize}
    \item ¿Cuál es la diferencia de la calidad del algoritmo propuesto en comparación a las versiones originales de los algoritmos de optimización implementados en la resolución de funciones benchmark y un problema real para un mismo contexto?. %esta pregunta nos permite averiguar qué tan bueno es el algoritmo dada una instancia en particular

    \item ¿Cuál es el impacto de la escalabilidad en la calidad de las soluciones del algoritmo con la mejora propuesta respecto al algoritmo en su versión original?, es decir, en instancias más grandes del problema. %esta pregunta nos permite saber si la instancia funciona mejor cuando la cantidad de puntos de entrega del problema es pequeño o grande

\end{itemize}










\section{Hipótesis}
\label{sc:HIPO}

%Define hipótesis nula y alternativa consistentes con la pregunta de investigación y las mediciones a realizar.%
Para esta investigación se planean hipótesis según el diseño experimental que se implementa. Debido a que la comparación es pareada entre el algoritmo con la propuesta y en su versión original, se utilizarán dos hipótesis para cada par de algoritmos que se utilice en este trabajo, las cuales son:
\begin{itemize}
    \item Hipótesis nula (H0): ``No hay diferencia significativa de mejora en el desempeño del algoritmo propuesto, para los problemas de optimización''.

    \item Hipótesis alternativa (H1): ``Existe una diferencia significativa de mejora en el desempeño de la técnica propuesta, en el proceso de resolución de problemas de optimización, comparado con el enfoque original del método que no incorporan la optimización del proceso de búsqueda.''.

\end{itemize}


Bajo esta hipótesis H0, se considera la posibilidad de que la propuesta no muestre una diferencia de mejora sobre el algoritmo originario comparado. Esta además implicaría, que cualquier diferencia en la calidad de soluciones entre el algoritmo propuesto y su versión de origen, no se debe a el algoritmo propuesto siendo superior, sino que podría deberse al azar o a otros factores. Por otro lado, en H1 se estaría asumiendo que la propuesta es superior en términos de calidad de solución comparado con el otro algoritmo de referencia.

El par de hipótesis para todo el diseño experimental tiene la finalidad de averiguar a nivel general si: ``Hay diferencia de mejora en la calidad del algoritmo propuesto para los problemas planteados en los experimentos comparado con todas las versiones originales del algoritmo en cuestión''.




\section{Objetivo General}
\label{sc:OG}
Evaluar el impacto del método bioinspirado mejorado, que incorpora técnicas de XAI y estrategias para la detección de estancamientos en óptimos locales, en la calidad de las soluciones obtenidas en problemas de optimización continua y discreta, comparado con la versión original del método, con el fin de determinar si la mejora en la optimización del proceso de búsqueda resulta en un desempeño significativamente superior.



\section{Objetivos Específicos}
\label{sc:OE}

\begin{itemize}
    \item Investigar métodos de explicabilidad de modelos XAI que permitan seleccionar las variables de decisión más relevantes del análisis bayesiano durante el proceso de optimización.
    \item Desarrollar un procedimiento para recopilar y calcular métricas de los datos generados por los algoritmos metaheurísticos durante el proceso de optimización.
    \item Validar la técnica propuesta utilizando problemas de optimización continua CEC LSGO 2023 y CEC 2017, junto a un problema discreta Profitability Tour Problem (PTP).
    \item Evaluar la eficacia de la técnica desarrollada en términos de la calidad de las soluciones obtenidas durante el proceso de optimización, comparándolas con enfoques tradicionales que no incorporan la técnica propuesta, utilizando métricas de rendimiento como el fitness, convergencia, calidad de solución, robustez (análisis estadístico).
    \item Reportar los resultados obtenidos de la propuesta, proporcionando una explicación detallada del comportamiento del algoritmo durante el proceso de optimización.
\end{itemize}

\chapter{Metodología}
   \label{sc:METOD}

Para este trabajo, se utilizará una metodología investigativa cuantitativa con enfoque experimental \cite{metodologia_investigacion} basada en \cite{10.5555/1718024}, la cual es ampliamente utilizada para analizar el comportamiento de las metaheurísticas. Esta metodología, se basa en 3 fases, la primera es el diseño experimental, la segunda las mediciones que serán realizadas en el experimento y por último definir cómo se hará el reporte de resultados y su visualización. En en la Figura \ref{fig:metodologia} se ve representada.


\begin{figure}[H]
    \centering
    \includegraphics[width=0.75\textwidth]{imagenes/disExpUpgrade.png}
    \caption{Metodología investigativa experimental basada en \cite{10.5555/1718024}}
    \label{fig:metodologia}
\end{figure}



\begin{enumerate}
    \item \textbf{Diseño experimental}: En esta fase, se deben definir claramente:
    \begin{itemize}
        \item \textbf{Definición de metas}: se deben definir claramente los objetivos de los experimentos. Estos objetivos pueden variar ampliamente y deben especificar qué se espera lograr o demostrar con el análisis de la metaheurística. Los objetivos ya fueron declarados en la sección anterior, tanto el general como los específicos.

        \item \textbf{Selección de instancias}: es crucial seleccionar adecuadamente las instancias de prueba (problemas de optimización combinatorial), asegurando que sean diversas en términos de tamaño, dificultad y estructura. Estas instancias se deben dividir en dos subconjuntos: uno para calibrar los parámetros de la metaheurística y otro para evaluar su rendimiento. La calibración de parámetros debe hacerse con cuidado para evitar el sobreajuste, que podría reducir la robustez de la metaheurística en instancias no conocidas.

        Para el presente trabajo, la calibración será omitida y solo se utilizarán los parámetros recomendados para la metaheurística en específico, que posteriormente serán plasmados en la sección de ``Configuración experimental de parámetros de metaheurísticas''.

        %Metric: Variables/Escalas/Tipos de Mediciones%
        También existen otro tipo de parámetros para los algoritmos de optimización, que se conocen como factores que pueden ser considerados generales para todas las metaheurísticas. Estas dos variables independientes son:
        %variables independientes
        \begin{itemize}
            \item número máximo de iteraciones
            \item número de individuos/soluciones/agentes por cada generación (también conocido como ``tamaño de la población'')
        \end{itemize}

        Por otra parte, para evaluar el rendimiento de estos algoritmos, se seleccionarán instancias llamadas ``funciones benchmark'', en particular las de la CEC 2013 y 2017 Large-Scale Global Optimization (LSGO) \cite{li2013benchmark, wu2017problem}. Por último, se probará la propuesta en una instancia de un problema de ingeniería de la vida real llamado, \textit{Profitable Tour Problem} (PTP) \cite{Vansteenwegen2019}.

        El PTP o Problema del Recorrido Rentable, es una variante del problema clásico del vendedor viajero (TSP \cite{LAPORTE1992231}, por sus siglas en inglés). En el PTP, tanto el beneficio como el costo de viaje se consideran en la función objetivo, convirtiendose en una función multiobjetivo. El objetivo principal es maximizar el beneficio total recolectado menos el costo total del viaje. Esto implica que no necesariamente se deben visitar todos los clientes, sino seleccionar aquellos cuya visita sea más rentable considerando la relación entre el beneficio obtenido y el costo de incluirlos en la ruta (considerado como unas de las diferencias con el TSP).

        Este problema es crucial e importante en cuanto a los beneficios que entrega tanto para las empresas de transporte como para los clientes. Para las empresas, el PTP optimiza las rutas al maximizar el beneficio neto, considerando costos y ganancias de visitar diferentes clientes. Esto mejora la eficiencia operativa y la rentabilidad, reduciendo costos y aumentando ingresos. La capacidad de determinar las rutas más rentables permite a las empresas utilizar mejor sus recursos, minimizar tiempos de inactividad y ofrecer un servicio más competitivo en un mercado exigente. En el ámbito de los clientes, el PTP permite ofrecer un servicio más eficiente y personalizado, priorizando entregas críticas y rentables. Esto mejora la satisfacción del cliente al recibir un servicio más rápido y adaptado a sus necesidades. La optimización de rutas basada en el PTP asegura que las entregas sean puntuales y confiables, lo que fortalece la relación con los clientes y mejora la reputación de la empresa. En conjunto, el PTP podría proporcionar una ventaja competitiva significativa, equilibrando costos y beneficios, y mejorando tanto la gestión de recursos como la adaptabilidad al mercado.

        \textbf{Definición matemática del PTP}: El PTP se puede formular como un modelo de programación entera con las siguientes variables de decisión:
        \begin{itemize}
            \item $x_{ij} = 1$ si se visita el nodo $i$ seguido del nodo $j$, y 0 en caso contrario.
            \item $v_i$ se usa en las restricciones de eliminación de subtours para determinar la posición de los nodos visitados en la ruta.
        \end{itemize}

        \textbf{Función objetivo}:
        \[
            \text{Maximizar} \sum_{i=1}^{|N|-1} \sum_{j=2}^{|N|} (P_i x_{ij} - t_{ij} x_{ij})
        \]
        Donde $P_i$ es el beneficio de visitar el nodo $i$ y $t_{ij}$ es el costo de viajar del nodo $i$ al nodo $j$.

        \textbf{Restricciones}:
        \begin{enumerate}
            \item La ruta comienza en el nodo 1 y termina en el nodo $|N|$:
            \[
            \sum_{j=2}^{|N|} x_{1j} = \sum_{i=1}^{|N|-1} x_{i|N|} = 1
            \]
            \item Cada nodo $k$ es visitado a lo sumo una vez:
            \[
            \sum_{i=1}^{|N|-1} x_{ik} = \sum_{j=2}^{|N|} x_{kj} \leq 1, \forall k = 2, \ldots, (|N|-1)
            \]
            \item Restricciones de eliminación de subtours basadas en la formulación de Miller–Tucker–Zemlin:
            \[
            2 \leq u_i \leq |N|, \forall i = 2, \ldots, |N|
            \]
            \[
            u_i - u_j + 1 \leq (|N| - 1)(1 - x_{ij}), \forall i, j = 2, \ldots, |N|
            \]
        \end{enumerate}

        \textbf{Eliminación de Subtours}: En muchos problemas de enrutamiento, formulaciones matemáticas incompletas pueden permitir la aparición de subtours en una solución. Las restricciones de subtour eliminan estas soluciones no válidas. Sin las restricciones de (c), la combinación de rutas como 1-5-1 y 2-3-4-2 sería aceptable, pero no es válida en términos de la definición del problema. Las restricciones para eliminar subtours en el PTP son las siguientes:

        \[
            2 \leq u_i \leq |N|, \forall i = 2, \ldots, |N|
        \]
        \[
            u_i - u_j + 1 \leq (|N| - 1)(1 - x_{ij}), \forall i, j = 2, \ldots, |N|
        \]

        Para cualquier posible subtour, estas restricciones se violarían, asegurando que no se formen subtours en la solución óptima \cite{Vansteenwegen2019}.





    \end{itemize}


    \item \textbf{Mediciones}: En esta fase se seleccionan las medidas de rendimiento y los indicadores que se van a calcular durante el proceso de optimización. Estas medidas o métricas son variables dependientes que pueden incluir:

    \begin{itemize}
        \item \textbf{Fitness}: cuantitativamente esta variable es ordinal y dependiente del contexto. Esta variable es el valor resultante de la función objetivo que se da en la ecuación perteniente a la instancia que se esta resolviendo. %permite evaluar la mejora de la solución que entrega el algoritmo --> tributa a la pregunta de invest. 1 y 3%

        \item \textbf{Tiempo de resolución o tasa de convergencia}: cuantitativamente esta variable es ordinal y dependiente del contexto. Es una medida de cuánto mejora el valor fitness por cada iteración. Una tasa de convergencia alta indica que el algoritmo está convergiendo rápidamente hacia la solución óptima. Esta se calcula de la siguiente manera:
        \begin{equation}
            TC = \frac{f(x_0) - f^*}{\Delta t}
        \end{equation}
        donde $f(x_0)$ es el valor fitness de la iteración inicial (0), $f^*$ es el valor fitness de la última iteración del algoritmo y $\Delta t$ es la cantidad máxima de iteraciones.
        %permite evaluar la rapidez con la que el algoritmo converge al optimo encontrado --> tributa a la pregunta de invest. 1 y 3%


        \item \textbf{Calidad de la solución:} Los indicadores de desempeño para definir la calidad de la solución en términos de precisión, se basa en medir la distancia (o el porcentaje) de desviación de la solución obtenida con respecto a la solución óptima del problema o una conocida hasta el momento. Cuantitativamente es una variable ordinal y dependiente del contexto. La medida calcula la diferencia que existe entre la solución conocida y la solución que se obtuvo en la última iteración del algoritmo en cuestión. Mientras más se acerque a 0, menos diferencia existe entre la mejor solución y la encontrada. La ecuación de esta es:
        \begin{equation}
            Calidad = \frac{f(x^*) - f_{optimo}}{max(|f(x*)|,|f_{optimo}|)}
        \end{equation}
        donde $f(x^*)$ es la solución encontrada por el algoritmo, y $f_{optimo}$ es el valor fitness conocido de la instancia/problema.


        \item \textbf{Robustez}: En general, la robustez se define como el nivel de insensibilidad que presenta el algoritmo, frente a pequeñas desviaciones en las entradas de las instancias, o en los parámetros de la metaheurística. Cuanto menor sea la variabilidad de las soluciones obtenidas, mejor será la robustez.

    \end{itemize}


    En adición a esto, en esta fase se detalla el diseño de todo el análisis de las variables dependientes resultantes de los experimentos (métricas anteriores) con el fin de ser interpretadas y poder dar una conclusión clara del estudio realizado. Se procederá definiendo un \textit{Análisis de datos atípicos}, el \textit{Análisis ordinario} y luego el \textit{Análisis estadístico}.


    %Definición de análisis de datos: Describe el test estadístico a utilizar y las variables sobre las cuales lo va a aplicar.
    \begin{itemize}
        \item \textbf{Análisis de datos atípicos}:en esta análisis se detectan datos atípicos para proceder a realizar una reducción de datos. Este es un tipo de análisis que se realiza sobre un conjunto de datos para determinar cuáles de estos son considerados como datos atípicos (también llamados en inglés \textit{outliers}) \cite{VijayBala2019}. Estos datos son valores extremos o lejanos al conjunto de datos de una muestra que anormalmente se encuentran fuera del patrón general de una distribución de variables \cite{kjae2017}. Como en la mayoría de las pruebas estadísticas asumen que los datos se distribuyen normalmente, la identificación de valores atípicos debe preceder al análisis de datos. Se utilizan diferentes métodos para identificar valores atípicos en una distribución normal \cite{kjae2017}. Al saber cuales son los valores que desvirtúan la muestra se clasifican por medio de dos tipos de datos outliers: \textit{Leves} y \textit{Extremos}. Un valor de muestra pertenecerá a una, u otra clasificación, dependiendo de su ubicación dentro de los rangos intercuartílicos.

        \begin{itemize}
            \item \textbf{Leves}: Sea $VM$ el valor de la muestra, $Q1$ el cuartil 1, $Q3$ el cuartil 3, y $IQR$ el rango inter-cuartil, se dice que $VM$ es outlier leve si es:
            \begin{align*}
                VM < (Q1 - 1.5 \cdot IQR) \\
                VM > (Q3 + 1.5 \cdot IQR)
            \end{align*}

            \item \textbf{Extremos}: Sea $VM$ el valor de la muestra, $Q1$ el cuartil 1, $Q3$ el cuartil 3, y $IQR$ el rango inter-cuartil, se dice que $VM$ es outlier extremo si es:
            \begin{align*}
                VM < (Q1 - 3 \cdot IQR) \\
                VM > (Q3 + 3 \cdot IQR)
            \end{align*}

            Luego de que los valores atípicos sean reconocidos, estos se remueven para evitar que la muestra se desvirtúe, y posteriormente se realizan nuevas pruebas para completar las 30 salidas por ejecución.

        \end{itemize}


        \item \textbf{Análisis de ordinario}: en este análisis se procede a buscar una serie de medidas estandarizadas que nos ayuden a contrastar los resultados obtenidos. Las medidas utilizadas son las siguientes:

        \begin{enumerate}
            \item \textbf{Media aritmética}: esta se define como el promedio aritmético de una distribución. Esta es la suma de todos los valores dividida entre el número de casos. Es una medida solamente aplicable a mediciones por intervalos o de razón \cite{metodologia_investigacion}. Esta medida se basa en el principio de la \textit{esperanza matemática} (valor esperado). La fórmula de esta es:
            \begin{align*}
                \overline{x} = \frac{\sum_{i=1}^{N} x_i}{N}
            \end{align*}
            donde $N$ es la cantidad total de valores a trabajar. $\overline{x}$ es la media del conjunto de $N$ valores, y $\sum_{i=1}^{N} x_i$ manifiesta la suma total desde el elemento $i$ hasta el $N$ de dicho conjunto.

            \item \textbf{Mejor valor}: es aquel valor, del total de valores obtenidos, que es superior al resto basándose en los requerimientos del problema. El mejor valor del conjunto puede ser muy grande (maximización), o bien muy pequeño (minimización).

            \item \textbf{Peor valor}: es aquel valor, del total de valores obtenidos, que es inferior al resto basándose en los requerimientos del problema. El peor valor del conjunto puede ser muy grande (minimización), o bien muy pequeño (maximización).

            \item \textbf{Desviación estándar}: es el promedio de desviación de las puntuaciones con respecto a la media. Esta medida se expresa en las unidades originales de medición de la distribución. Se interpreta en relación con la media. Cuanto mayor sea la dispersión de los datos alrededor de la media, mayor será la desviación estándar \cite{metodologia_investigacion}. Por lo que, esta mide la distancia a la que se encuentran cada uno de los valores de un conjunto respecto a la media aritmética de este. Esta se simboliza con una $s$ o $\sigma$ y se calcula de la siguiente manera:
            \begin{align*}
                \sigma = \sqrt{ \frac{ \sum_{i=1}^{N} (x_i - \overline{x})^2 }{N - 1} }
            \end{align*}

            \item \textbf{Rango intercuartílico}: este es el rango de valores que incluye el 50\% medio de una distribución \cite{Kipfer2021}. Denotado por las siglas \textit{IQR} es una medida de dispersión estadística que se obtiene de la \textit{diferencia entre el tercer y primer cuartil} de una distribución y que además sirve para encontrar valores atípicos \cite{10.1007/978-981-10-7563-6_53}. Las definiciones de primer y tercer cuartil, junto a sus respectivas fórmulas para calcularlos, se muestran a continuación.

            \begin{itemize}
                \item \textbf{Primer cuartil}: es aquel valor límite que nos indica que el 25\% de los datos de la distribución estadística son menores o iguales a él. Para determinar dicho valor se utiliza la siguiente fórmula:
                \begin{align*}
                    Q_1 = \frac{N+1}{4}
                \end{align*}

                \item \textbf{Tercer cuartil}: es aquel valor límite que nos indica que el 75\% de los datos de la distribución estadística son menores o iguales a él. Para determinar dicho valor se utiliza la siguiente fórmula:
                \begin{align*}
                    Q_3 = \frac{3(N+1)}{4}
                \end{align*}

            \end{itemize}

        \end{enumerate}


        \textbf{Análisis estadístico}: la estadística es la metodología que los científicos y matemáticos han desarrollado para interpretar y sacar conclusiones de los datos \cite{Anderson1996}.

        Un análisis estadístico debe ser enfocado de diversas \textit{pruebas estadísticas}. Una prueba estadística es, en palabras sencillas, una forma de evaluar la evidencia entregada por un conjunto de datos disponibles para respaldar una hipótesis propuesta. Los pasos necesarios para realizar una prueba estadística son:

        \begin{enumerate}
            \item \textbf{Hipótesis}: en este se debe plantear la \textit{Hipótesis nula} ($H_0$), que es la afirmación de que una serie de parámetros (o fenómenos) no tienen relación entre sí. Se utiliza como punto de partida en cualquier investigación, y se rechaza solamente si los datos de la muestra evidencian lo contrario a lo planteado. Por otro lado, la \textit{Hipótesis investigativa} ($H_1$) es aquello que el investigador especula que sucederá al realizar un experimento, al finalizar una investigación, o tras ocurrir un cierto acontecimiento. Para poder aceptar $H_1$ se debe rechazar $H_0$ \cite{metodologia_investigacion}.

            \item \textbf{Nivel de significancia}: representado por la letra griega $\alpha$, es un valor umbral que nos indica si el resultado de un estudio es estadísticamente significante \cite{Devore2012}. Un resultado se dice que es estadísticamente significante cuando es improbable que este surgiera por obra del azar.

            El nivel de significancia representa la probabilidad de aceptar la hipótesis nula, y por lo general en metaheurísticas se suele limitar a 0.05 (5\%). En otras palabras, ese cinco por ciento representa a aquellos casos contrario a $H_1$ (favorables a $H_0$) que, se asume, surgieron por obra del azar.

            \item \textbf{Valor calculado}: También conocido como \textit{Estadístico de Prueba}, es una variable aleatoria que mide el grado de concordancia entre una muestra de datos y la hipótesis nula \cite{Devore2012}. En otras palabras, el estadístico de prueba compara los datos de una muestra con los resultados esperados bajo la hipótesis nula, con el objetivo de poder determinar si es posible o no rechazar dicha hipótesis. Este estadístico se utiliza para calcular el \textit{Valor P} y, dependiendo de la prueba a realizar, el estadístico empleado cambia. Las distintas pruebas y estadísticos pueden observarse en la Tabla \ref{tab:pruebas-hipo}.

            \begin{table}[ht!]
            \centering
            \begin{tabular}{|c|c|}\hline
                Prueba de hipótesis & Estadístico de prueba \\\hline

                Pruebas Z & Estadísticos Z \\

                Pruebas t & Estadísticos t \\

                Anova & Estadísticos F \\

                Pruebas chi-cuadrada & Estadístico chi-cuadrada \\

                Kolmogorov-Smirnov & Máxima diferencia \\

                Shapiro-Wilk & Estadístico W \\\hline
            \end{tabular}
            \caption{Prueba de hipótesis y estadísticos de prueba.}
            \label{tab:pruebas-hipo}
            \end{table}

            \item \textbf{Valor P}: es la probabilidad de cometer un \textit{Error tipo I}. Un error tipo I es cuando no se rechaza $H_1$ a pesar de que el evento ocurrido (proposición) no cabe dentro de dicha hipótesis \cite{Iversen1997}. En otras palabras, un error tipo I es cuando ocurre un evento contrario a lo estadísticamente esperado.

            \item \textbf{Decisión}: es la decisión de rechazar o no rechazar la hipótesis nula ($H_0$). Para saber cuando hacer esto, $H_0$ se debe tener en cuenta lo siguiente:
            \begin{itemize}
                \item \textbf{Rechazo}: si se rechaza $H_0$, esto implica no rechazar $H_1$. $H_0$ se rechaza si $P \leq \alpha$.

                \item \textbf{No se rechaza}: si se no se rechaza $H_0$, esto implica rechazar $H_1$. $H_0$ no se rechaza si $P > \alpha$.

            \end{itemize}

            \item \textbf{Conclusión}: Una vez realizados todos los pasos anteriores se debe llegar a una conclusión. La conclusión debe ir acorde a las hipótesis planteadas ($H_0$ y $H_1$) y, básicamente, sirve como reafirmación final sobre la hipótesis aceptada.

        \end{enumerate}


        \begin{figure}[ht!]
            \centering
            \includegraphics[width=0.9\textwidth]{imagenes/analisis_estadistico.png}
            \caption{Análisis estadístico junto a sus ramas y pruebas estadísticas.}
            \label{fig:analisisEsta}
        \end{figure}

        Ya explicado el funcionamiento de una prueba estadística, se detallan los test mencionados en la Figura \ref{fig:analisisEsta}. Para este estudio se utilizará lo que esta en marcado en azul en la imagen. Debido a la cantidad de pruebas y datos a analizar no superarán las 30 muestras, se utilizará las pruebas de estadístico no-paramétrico de Shapiro-Wilk. También se utilizará Wilcoxon Mann Whitney, ya que las muestras de los algoritmos no provienen de la naturaleza (es decir, que no debieran tener una distribución normal) y porque cada ejecución es totalmente independiente a la anterior. De acuerdo a estas pruebas, a continuación se indicará cuándo utilizarlas, para qué sirven y cómo desarrollarlas.

        \begin{itemize}
            \item \textbf{Prueba de Shapiro-Wilk}: test utilizado cuando el tamaño de la muestra es menor a 50 datos ($N < 50$). El propósito de este test es contrastar la normalidad (distribución) de un conjunto de datos \cite{Hanusz_Tarasinska_Zielinski_2016}. Para utilizar este test se debe partir de la premisa ($H_0$) de que una muestra $\{x_1, x_2, ..., x_n\}$ proviene de una población normalmente distribuida. Para saber si dicha hipótesis $H_0$ es rechazada o no rechazada, primero se debe ordenar la muestra de menor a mayor, y luego utilizar el siguiente estadístico de contraste:

            \begin{gather*}  \label{eq:shapiro_wilk}
                W = \frac{1}{n \cdot s^2} ( \sum_{j=1}^{h} a_{ni} (x_{n-i+1} - x_i))^2 \\
                \wedge \;
                f(n) =  \begin{cases}
                    \frac{n}{2},  & \text{ si } n \text{ es par} \\
                    \frac{n-1}{2}, & \text{ si } n \text{ es impar} \\
                \end{cases}
            \end{gather*}

            donde $x_j$ es el j-ésimo valor muestral del ordenamiento, $s_2$ es la varianza muestral, $a_{ni}$ es un coeficiente tabulado, y $h$ es una variable real que puede tomar dos valores dependiendo del tamano de la muestra, mostrados en \ref{eq:shapiro_wilk}. Obtenido el valor de $W$, su distribución permite calcular el valor crítico del test, del cual se toma la decisión sobre factibilidad de la normalidad de la muestra.

            \item \textbf{Prueba de Wilcoxon Mann Whitney}: también conocido como \textit{Test U}, es una prueba no paramétrica (los datos no se ajustan a una distribución conocida) que se aplica sobre dos muestras independientes \cite{fagerland2009wilcoxon}. El objetivo de esta prueba es comprobar la similitud entre ambas muestras a partir de las diferencias que puedan detectarse en ellas. Para utilizar este test se debe partir de la premisa ($H_0$) de que las distribuciones de ambas muestras son las mismas. Para saber si dicha hipótesis es rechazada o no rechazada, primero se debe ordenar ambas muestras de menor a mayor, y luego utilizar el siguiente estadístico de contraste:

            \begin{subequations} \label{eq:wmw}
                \begin{equation}
                    U = \sum_{i=1}^{m} R(Y_i)
                \end{equation}
                \begin{equation}
                    T = \sum_{i=1}^{n} R(X_i)
                \end{equation}
                \begin{equation}
                    U + T = \frac{N(N+1)}{2}
                \end{equation}
            \end{subequations}

            donde $m$ y $n$ son los tamaños de la primera y segunda muestra respectivamente, y $R(Y_i)$ junto a $R(X_i)$ son las funciones de rangos de dichas muestras. Una vez obtenidos los valores $U$ y $T$, se comprueba que se cumpla la Ec. \ref{eq:wmw}.

        \end{itemize}


    \end{itemize}


    \item \textbf{Reporte}: la interpretación de los resultados es de forma crítica basada en los datos obtenidos de los experimentos realizados. Para efectuar lo anterior la muestra de resultados y el análisis se realiza mediante tablas de comparación, herramientas gráficas, como barras de desviación (intervalos de confianza, diagramas de caja o violín) y gráficas de convergencia hacia el óptimo encontrado en la ejecución. Esto permitirá tener una mejor comprensión de la evaluación del desempeño de los resultados obtenidos. La interpretación de los resultados debe estar alineada con los objetivos definidos inicialmente, y se debe asegurar la reproducibilidad de los experimentos computacionales, proporcionando suficiente detalle sobre los procedimientos seguidos y los parámetros utilizados.

\end{enumerate}

\chapter{Planificación}
   
% --------------------------------------------------------------------------

\section{Plan de trabajo}
\label{sc:PT}
La planificación esta desglosada en 2 paquetes de trabajo (WP). El primero (WP$_{1}$) corresponde a la implementación de la técnica híbrida en las metaheurísticas seleccionadas con pruebas preliminares de su funcionamiento. Luego, el segundo paquete (WP$_{2}$) tratará la fase de experimentación junto al análisis de resultados, con la posterior interpretación de estos y conclusión. El trabajo de paquetes (WP) están separados en actividades (A$_{x.y}$) para presentar el plan de trabajo mejor (ver la Tabla~\ref{tbl:wp}).



\begin{comment}

\begin{table}[!ht]
\begin{center}
\begin{scriptsize}
\caption{Gráfico Gantt - Semestre 2, 2024}\label{tbl:wp}
\vspace{5pt}
\begin{tabular}{|p{.42cm}|p{.15cm}|p{.15cm}|p{.15cm}|p{.15cm}|p{.15cm}|p{.15cm}|p{.15cm}|p{.15cm}|p{.15cm}|p{.15cm}|p{.15cm}|p{.15cm}|p{.15cm}|p{.15cm}|p{.15cm}|p{.15cm}|p{.15cm}|p{.15cm}|p{.15cm}|p{.15cm}|}
\hline
\rowcolor{gray!20}
\multicolumn{1}{|c|}{}
& \multicolumn{4}{c|}{\textbf{AGOSTO}}
& \multicolumn{4}{c|}{\textbf{SEPTIEMBRE}}
& \multicolumn{4}{c|}{\textbf{OCTUBRE}}
& \multicolumn{4}{c|}{\textbf{NOVIEMBRE}}
& \multicolumn{4}{c|}{\textbf{DICIEMBRE}} \\ \cline{2-21}
\rowcolor{gray!20}
\multicolumn{1}{|c|}{\multirow{-2}{*}{\textbf{Act.}}} &
\multicolumn{1}{c|}{\textbf{1}} &
\multicolumn{1}{c|}{\textbf{2}} &
\multicolumn{1}{c|}{\textbf{3}} &
\multicolumn{1}{c|}{\textbf{4}} &
\multicolumn{1}{c|}{\textbf{1}} &
\multicolumn{1}{c|}{\textbf{2}} &
\multicolumn{1}{c|}{\textbf{3}} &
\multicolumn{1}{c|}{\textbf{4}} &
\multicolumn{1}{c|}{\textbf{1}} &
\multicolumn{1}{c|}{\textbf{2}} &
\multicolumn{1}{c|}{\textbf{3}} &
\multicolumn{1}{c|}{\textbf{4}} &
\multicolumn{1}{c|}{\textbf{1}} &
\multicolumn{1}{c|}{\textbf{2}} &
\multicolumn{1}{c|}{\textbf{3}} &
\multicolumn{1}{c|}{\textbf{4}} &
\multicolumn{1}{c|}{\textbf{1}} &
\multicolumn{1}{c|}{\textbf{2}} &
\multicolumn{1}{c|}{\textbf{3}} &
\multicolumn{1}{c|}{\textbf{4}} \\ \hline
\multicolumn{1}{c}{\textbf{WP}$_{1}$} \\ \hline
A$_{1.1}$ &X&X&X&X& & & & & & & & & & & & & & & & \\ \hline
A$_{1.2}$ & & &X&X&X&X&X& & & & & & & & & & & & & \\ \hline
A$_{1.3}$ & & &X&X& & & & & & & & & & & & & & & & \\ \hline
A$_{1.4}$ & & & & & & & &X&X&X&X&X&X& & & & & & & \\ \hline
A$_{1.5}$ & & & & & & & & & & & &X&X&X& & & & & & \\ \hline
\multicolumn{1}{c}{\textbf{WP}$_{2}$} \\ \hline
A$_{2.1}$ & & & & & & & & & & & & & &X&X& & & & & \\ \hline
A$_{2.2}$ & & & & & & & & & & & & & & &X&X& & & & \\ \hline
A$_{2.3}$ & & & & & & & & & & & & & & & & &X&X& & \\ \hline
A$_{2.4}$ & & & & & & & & & & & & & & & & & &X&X&X \\ \hline
\end{tabular}
\end{scriptsize}
\end{center}
\end{table}

\end{comment}


\begin{table}[!ht]
\begin{center}
\begin{scriptsize}
\caption{Gráfico Gantt - Semestre 2, 2024}\label{tbl:wp}
\vspace{5pt}
\begin{tabular}{|p{.42cm}|p{.6cm}|p{.6cm}|p{.6cm}|p{.6cm}|p{.6cm}|p{.6cm}|p{.6cm}|p{.6cm}|p{.6cm}|p{.6cm}|p{.6cm}|p{.6cm}|p{.6cm}|p{.6cm}|p{.6cm}|p{.6cm}|p{.6cm}|p{.6cm}|p{.6cm}|p{.6cm}|}
\hline
\rowcolor{gray!20}
\multicolumn{1}{|c|}{}
& \multicolumn{4}{c|}{\textbf{AGOSTO}}
& \multicolumn{4}{c|}{\textbf{SEPTIEMBRE}}
& \multicolumn{4}{c|}{\textbf{OCTUBRE}}
& \multicolumn{4}{c|}{\textbf{NOVIEMBRE}}
& \multicolumn{4}{c|}{\textbf{DICIEMBRE}} \\ \cline{2-21}
\rowcolor{gray!20}
\multicolumn{1}{|c|}{\multirow{-2}{*}{\textbf{Act.}}} &
\multicolumn{1}{c|}{\textbf{1}} &
\multicolumn{1}{c|}{\textbf{2}} &
\multicolumn{1}{c|}{\textbf{3}} &
\multicolumn{1}{c|}{\textbf{4}} &
\multicolumn{1}{c|}{\textbf{1}} &
\multicolumn{1}{c|}{\textbf{2}} &
\multicolumn{1}{c|}{\textbf{3}} &
\multicolumn{1}{c|}{\textbf{4}} &
\multicolumn{1}{c|}{\textbf{1}} &
\multicolumn{1}{c|}{\textbf{2}} &
\multicolumn{1}{c|}{\textbf{3}} &
\multicolumn{1}{c|}{\textbf{4}} &
\multicolumn{1}{c|}{\textbf{1}} &
\multicolumn{1}{c|}{\textbf{2}} &
\multicolumn{1}{c|}{\textbf{3}} &
\multicolumn{1}{c|}{\textbf{4}} &
\multicolumn{1}{c|}{\textbf{1}} &
\multicolumn{1}{c|}{\textbf{2}} &
\multicolumn{1}{c|}{\textbf{3}} &
\multicolumn{1}{c|}{\textbf{4}} \\ \hline
\multicolumn{1}{c}{\textbf{WP}$_{1}$} \\ \hline
A$_{1.1}$ &X&X&X&X& & & & & & & & & & & & & & & & \\ \hline
A$_{1.2}$ & & &X&X&X&X&X& & & & & & & & & & & & & \\ \hline
A$_{1.3}$ & & &X&X& & & & & & & & & & & & & & & & \\ \hline
A$_{1.4}$ & & & & & & & &X&X&X&X&X&X& & & & & & & \\ \hline
A$_{1.5}$ & & & & & & & & & & & &X&X&X& & & & & & \\ \hline
\multicolumn{1}{c}{\textbf{WP}$_{2}$} \\ \hline
A$_{2.1}$ & & & & & & & & & & & & & &X&X& & & & & \\ \hline
A$_{2.2}$ & & & & & & & & & & & & & & &X&X& & & & \\ \hline
A$_{2.3}$ & & & & & & & & & & & & & & & & &X&X& & \\ \hline
A$_{2.4}$ & & & & & & & & & & & & & & & & & &X&X&  \\ \hline
\end{tabular}
\end{scriptsize}
\end{center}
\end{table}




\begin{itemize}

    \item \textbf{WP$_1$: Implementación de la técnica híbrida en las metaheurísticas para los problemas de optimización y ajuste de parámetros.}

    \begin{itemize}

        \item \textbf{A$_{1.1}$} Implementación de al menos 4 metaheurísticas para resolver funciones benchmark.

        \item \textbf{A$_{1.2}$} Estudiar, diseñar, y documentar en código los procedimientos de la propuesta para integrarlas a la metaheurística.

        \item \textbf{A$_{1.3}$} Ajustar los parámetros de las metaheurísticas a través de las recomendaciones hechas en los trabajos originales del algoritmo de optimización.

        \item \textbf{A$_{1.4}$} Implementar las metaheurísticas adaptadas al problema de optimización Profitable Tour Problem.

        \item \textbf{A$_{1.5}$} Diseñar e implementar métodos de guardado de métricas mencionadas en la metodología del Capítulo 4 y funciones que almacenen los gráficos de estos datos.

    \end{itemize}

    \item \textbf{WP$_2$: Experimentación, análisis de resultados e interpretación}

    \begin{itemize}

        \item \textbf{A$_{2.1}$} Aplicar metodología de diseño experimental para unificar los datos almacenados post-ejecución, verificando su correctitud.

        \item \textbf{A$_{2.2}$} Ejecutar una larga fase de experimentación con los algoritmos con la propuesta incorporada y la versión original de estos en las instancias seleccionadas de los problemas de optimización combinatorial benchmarks y el Profitable Tour Problem. Conjuntamente ir verificando salidas para evitar errores.

        \item \textbf{A$_{2.3}$} Análisis estadístico comparativo entre el algoritmo con la nueva propuesta y su versión tradicional. Se evalúa la calidad de la solución y su rendimiento a través de la convergencia. Plasmados en tablas y gráficos.

        \item \textbf{A$_{2.4}$} Interpretación de los resultados, conclusiones, posibles consideraciones y trabajos a futuro.

    \end{itemize}

\end{itemize}




%%%%%%%%%% en el ámbito de la hibridación de metaheurísticas y métodos de inteligencia artificial, definición del problema, propuesta de solución, métodos investigativos para ejecutar el desarrollo experimental y la planificación

\chapter{Recursos}
   
\section{Recursos humanos}
\label{sc:RH}

\subsection{Equipo de Trabajo}

El equipo de trabajo necesario para este proyecto de investigación está compuesto por los siguientes roles y personal:

\begin{itemize}
    \item \textbf{Profesor Guía:} Responsable de guiar al investigador principal en sus dudas generales del acerca del tópico del proyecto de investigación.
    \begin{itemize}
        \item \textit{Nombre del profesor guía:} Dr. Rodrigo Olivares.
    \end{itemize}

    \item \textbf{Investigador Principal:} Responsable de idear la propuesta de solución al problema planteado de investigación aplicada, toma de decisiones estratégicas, planificación y desarrollo de todas las fases de la investigación.
    \begin{itemize}
        \item \textit{Nombre del investigador principal:} Sebastián M. Guzmán.
    \end{itemize}

\end{itemize}






\section{Recursos materiales}
\label{sc:RM}

\subsection{Hardware (HW)}
El hardware necesario para este trabajo de investigación incluye:
\begin{itemize}
    \item \textbf{Computador de Alto Rendimiento (HPC):} Equipadas con procesadores de última generación (mínimo Intel Core i7 o AMD Ryzen 7), 16GB de RAM y discos duros SSD de al menos 512GB.
    \item \textbf{Servidor:} Servidor dedicado al almacenamiento y procesamiento de datos con capacidades superiores al HPC (en caso de ser necesario).
\end{itemize}

\subsection{Software (SW)}
El software requerido incluye:
\begin{itemize}
    \item \textbf{Sistemas Operativos:} Windows 11 Home, Linux (Ubuntu 24.04 LTS) o macOS 14.
    \item \textbf{Software de desarrollo de código:} Python (con bibliotecas como NumPy, Pandas, SciPy, Sklearn o respositorios Github).
    \item \textbf{Herramientas de Desarrollo:} Entornos de desarrollo integrados (IDE) como Visual Studio Code, y Jupyter Notebook (o Colab de Google).
    \item \textbf{Software de Gestión de Proyectos:} Microsoft Project o Project Libre.
\end{itemize}

\subsection{Bibliotecas}
Para la consulta de literatura y referencias, se requiere acceso a:
\begin{itemize}
    \item \textbf{Biblioteca Digital:} Acceso las bases de datos académicas como ScienceDirect (Elsevier), SpringerLink, y Web of Science (WoS) para obtener libros electrónicos, artículos científicos, y publicaciones digitales relevantes al área de investigación.
\end{itemize}


%\chapter{Conclusiones preliminares}
%   
Las conclusiones no deben introducción nueva información además de la que ya se mencionó en el resto del documento. Deben ser presentadas en orden de importancia (más a menos).

Visión crítica. Presentar las posibles limitaciones y riesgos de su propuesta.

Indicar los trabajos futuros.



\appendix
%Este capítulo es solo instructivo: quitar o comentar
\include{gen}
%\include{appendix1}
%\include{appendix2}

%% REFERENCIAS BIBLIOGRÁFICAS: EN FORMATO IEEE, RECIENTES (ÚLTIMOS 10 AÑOS), EN INGLÉS Y ESPAÑOL.
\bibliographystyle{IEEEtran}
\bibliography{referencias}

\end{document}
